% Rough-drafting helpers.
% Toggle draft mode with \drafttrue or \draftfalse in ccpaper.tex.
% Toggle dyslexia mode with \dyslexictrue or \dyslexicfalse in ccpaper.tex.

\makeatletter
\@ifundefined{ifdraft}{\newif\ifdraft\drafttrue}{}
\@ifundefined{ifdyslexic}{\newif\ifdyslexic\dyslexicfalse}{}
\makeatother
\providecommand{\opendyslexicfontpath}{fonts/}

\usepackage{ifthen}
\usepackage{iftex}
\usepackage{lineno}
\usepackage{todonotes}
\usepackage[normalem]{ulem}

\ifdyslexic
  \ifPDFTeX
    \PackageWarning{custom_defs}{Dyslexia mode requires XeLaTeX or LuaLaTeX to use the installed OpenDyslexic font. Compile with xelatex or lualatex.}%
  \else
    \usepackage{fontspec}
    \IfFileExists{\opendyslexicfontpath OpenDyslexic-Regular.otf}{%
      \setmainfont{OpenDyslexic-Regular.otf}[
        Path=\opendyslexicfontpath,
        Ligatures=TeX,
        BoldFont=OpenDyslexic-Bold.otf,
        ItalicFont=OpenDyslexic-Italic.otf,
        BoldItalicFont=OpenDyslexic-BoldItalic.otf
      ]
      \setsansfont{OpenDyslexic-Regular.otf}[
        Path=\opendyslexicfontpath,
        Ligatures=TeX,
        BoldFont=OpenDyslexic-Bold.otf,
        ItalicFont=OpenDyslexic-Italic.otf,
        BoldItalicFont=OpenDyslexic-BoldItalic.otf
      ]
      \setmonofont{OpenDyslexicMono-Regular.otf}[
        Path=\opendyslexicfontpath,
        Ligatures=TeX
      ]
    }{%
      \PackageWarning{custom_defs}{Dyslexia mode requested, but OpenDyslexic-Regular.otf was not found at \opendyslexicfontpath. Install it from https://opendyslexic.org/ or redefine \string\opendyslexicfontpath.}%
    }
  \fi
\fi

\ifdraft
  \linenumbers
\fi

\newcommand{\draftonly}[1]{%
  \ifdraft #1\fi
}

\newcommand{\finalonly}[1]{%
  \ifdraft\else #1\fi
}

\newcommand{\needcite}{%
  \draftonly{\textcolor{red}{[citation needed]}}%
}

\newcommand{\placeholder}[1]{%
  \draftonly{\textcolor{gray}{[TODO: #1]}}%
}

\newcommand{\revise}[1]{%
  \draftonly{\textcolor{orange}{[REVISE: #1]}}%
}

\newcommand{\draftnote}[1]{%
  \draftonly{\todo[color=yellow!30,size=\small]{#1}}%
}

\newcommand{\inlinecomment}[2]{%
  \draftonly{\textcolor{#1}{\textbf{\textless\textless #2\textgreater\textgreater}}}%
}

\newcommand{\addtext}[1]{%
  \draftonly{\textcolor{blue}{#1}}\finalonly{#1}%
}

\newcommand{\removetext}[1]{%
  \draftonly{\textcolor{red}{\sout{#1}}}%
}

% Reviewer-style inline comments.
% Example: \reviewerone{Clarify this claim.}

\definecolor{revieweronecolor}{HTML}{C62828}
\definecolor{reviewertwocolor}{HTML}{1565C0}
\definecolor{reviewerthreecolor}{HTML}{6A1B9A}
\definecolor{reviewerfourcolor}{HTML}{00897B}
\definecolor{reviewerfivecolor}{HTML}{EF6C00}

\newcommand{\reviewercomment}[3]{%
  \draftonly{\textcolor{#1}{\textbf{\textless\textless REVIEWER \##2: #3\textgreater\textgreater}}}%
}

\newcommand{\reviewerone}[1]{\reviewercomment{revieweronecolor}{1}{#1}}
\newcommand{\reviewertwo}[1]{\reviewercomment{reviewertwocolor}{2}{#1}}
\newcommand{\reviewerthree}[1]{\reviewercomment{reviewerthreecolor}{3}{#1}}
\newcommand{\reviewerfour}[1]{\reviewercomment{reviewerfourcolor}{4}{#1}}
\newcommand{\reviewerfive}[1]{\reviewercomment{reviewerfivecolor}{5}{#1}}

% Figure and table snippets.
\newcommand{\figureplaceholder}[3]{%
  \begin{figure}[H]
    \centering
    \fbox{\parbox[c][0.25\textheight][c]{0.85\textwidth}{\centering #2}}
    \caption{#3}
    \label{fig:#1}
  \end{figure}
}

\newcommand{\twopanelfigureplaceholder}[3]{%
  \begin{figure}[H]
    \centering
    \begin{subfigure}{0.48\textwidth}
      \centering
      \fbox{\parbox[c][0.2\textheight][c]{0.9\textwidth}{\centering Panel A}}
      \caption{#1}
    \end{subfigure}
    \hfill
    \begin{subfigure}{0.48\textwidth}
      \centering
      \fbox{\parbox[c][0.2\textheight][c]{0.9\textwidth}{\centering Panel B}}
      \caption{#2}
    \end{subfigure}
    \caption{#3}
  \end{figure}
}

\newcommand{\landscapetableplaceholder}[1]{%
  \begin{landscape}
    \begin{table}[H]
      \centering
      \caption{#1}
      \begin{tabular}{lllll}
        \toprule
        Column 1 & Column 2 & Column 3 & Column 4 & Column 5 \\
        \midrule
        Value & Value & Value & Value & Value \\
        \bottomrule
      \end{tabular}
    \end{table}
  \end{landscape}
}

\newcommand{\basictableplaceholder}[1]{%
  \begin{table}[H]
    \centering
    \caption{#1}
    \begin{tabular}{lll}
      \toprule
      Column 1 & Column 2 & Column 3 \\
      \midrule
      Value & Value & Value \\
      \bottomrule
    \end{tabular}
  \end{table}
}
